% =====================================================================
\part{Conclusion et perspectives}
% =====================================================================

\chapter{Synthèse}

Ce cours a présenté les bases physiques de la thermique du bâtiment. La chaleur se transmet via les trois modes de transferts de chaleur qui sont la conduction, la convection et le rayonnement.
L'air est humide, il contient de la vapeur d'eau qui migre à travers les parois et peut condenser dès qu'elle atteint son point de rosée.

L'architecte n'a pas besoin de savoir résoudre les équations présentées dans ce cours. En revanche, il doit en comprendre les mécanismes pour faire des choix en phase de conception : implantation, orientation, volumétrie, composition des parois, dessin des détails. Ce sont ces choix architecturaux, qui déterminent le confort des usagers, la consommation d'énergie et l'impact carbone du projet. Les équipements (chauffage, ventilation, climatisation) viennent ensuite compléter une conception déjà performante (et non compenser des erreurs ou oublis de conception.

\chapter{Interventions}

\section{Projet neuf}

En projet neuf, l'architecte dispose de beaucoup de liberté. Voici les principales questions à se poser, à peu près dans l'ordre où elles apparaissent dans le projet :

\begin{enumerate}
    \item \textbf{Le site et le climat} : quel est le climat local ? D'où vient le soleil ? Quels sont les masques (bâtiments voisins, végétation, relief) ?
    \item \textbf{L'implantation et l'orientation} : orienter le bâtiment pour capter les apports solaires en hiver et pouvoir s'en protéger en été. Le choix de l'emplacement des espaces tampons
    \item \textbf{La volumétrie} : privilégier une volumétrie compacte, qui limite la surface d'enveloppe en contact avec l'extérieur, et donc les déperditions.
    \item \textbf{Les ouvertures} : dimensionner et répartir les surfaces vitrées selon l'orientation, et prévoir dès l'esquisse les protections solaires extérieures (casquettes, brise-soleil, volets, végétation).
    \item \textbf{L'enveloppe} : choisir une stratégie d'isolation continue et dessiner les détails pour traiter les ponts thermiques (jonctions mur/plancher, balcons, menuiseries). Assurer l'étanchéité à l'air.
    \item \textbf{La gestion de la vapeur} : ordonner les couches des parois pour que leur perméabilité à la vapeur augmente de l'intérieur vers l'extérieur, ou prévoir une membrane adaptée. Vérifier le comportement des parois en hiver et en été (avec Ubakus par exemple).
    \item \textbf{L'inertie et la ventilation} : conserver de la masse côté intérieur et permettre la ventilation naturelle traversante, notamment la nuit en été. Prévoir le renouvellement d'air et l'évacuation de l'humidité des pièces humides.
    \item \textbf{Les matériaux} : privilégier des matériaux à faible impact carbone (biosourcés, géosourcés, réemploi) en s'appuyant sur les FDES.
    \item \textbf{La réglementation} : garder en tête les six indicateurs de la RE 2020, en dialogue avec le bureau d'études thermiques. Le $B_{bio}$, le $I_c Construction$ et les degrés-heures d'inconfort dépendent directement des choix  de l'architecte.
\end{enumerate}

\section{Projet de réhabilitation}

En réhabilitation, le bâtiment existe déjà : son orientation, sa volumétrie et une partie de ses matériaux sont imposés par l'existant. La RE 2020 ne s'applique pas (pour le moment...), mais les enjeux de confort, d'énergie et de carbone restent les mêmes. Réhabiliter plutôt que démolir est d'ailleurs le choix le plus sobre en carbone. Voici les principales questions à se poser :

\begin{enumerate}
    \item \textbf{Le diagnostic} : avant toute intervention, comprendre comment le bâtiment fonctionne. De quoi sont faits les murs ? Comment gèrent-ils l'humidité ? Y a-t-il des désordres (humidité, moisissures, fissures) ? Où sont les principales déperditions ?
    \item \textbf{Les priorités} : intervenir d'abord là où les pertes sont les plus importantes. (La toiture, les murs,les menuiseries, le plancher bas.)
    \item \textbf{La stratégie d'isolation} : choisir entre isolation par l'extérieur (plus performante, conserve l'inertie, mais modifie la façade) et isolation par l'intérieur (préserve la façade, notamment patrimoniale, mais crée des ponts thermiques et réduit la surface habitable).
    \item \textbf{Le respect des murs anciens} : les murs anciens (pierre, brique, pisé, pan de bois) fonctionnent souvent grâce à leur capacité à laisser migrer la vapeur et à sécher. Les isoler avec des matériaux ou des membranes inadaptés peut piéger l'humidité et provoquer des désordres importants. Il s'agit de privilégier des matériaux perspirants (fibre de bois, chanvre, enduits à la chaux ou à la terre).
    \item \textbf{Les ponts thermiques} : identifier ceux que l'intervention ne pourra pas supprimer (planchers, refends, balcons) et vérifier qu'ils ne provoqueront pas de condensation de surface.
    \item \textbf{La ventilation} : une enveloppe rénovée est souvent plus étanche à l'air. Il peut être necessaire de prévoir un renouvellement d'air adapté, (pour éviter que l'humidité provoque de la condensation et moisissures).
    
\end{enumerate}

Dans tous les cas, une réhabilitation se pense de manière globale : intervenir sur un seul élément (changer les fenêtres sans ventiler, isoler sans traiter l'humidité) peut créer plus de problèmes qu'il n'en résout.
\chapter{Pour aller plus loin}
La bibliographie complète est présentée en fin de document, voici ici une liste plus courte de ressources utiles à consulter dans différents format :

\begin{encadre}{Livre}
  \begin{itemize}
    \item \textit{La Thermique du bâtiment}, O. Benjeddou (2016)
  \end{itemize}
\end{encadre}

\begin{encadre}{Site Web}
  \begin{itemize}
    \item AMACO plateforme ressources, \url{https://ressources.amaco.org/}
    \item UBAKUS, \url{https://www.ubakus.de/fr/calculateur-valeur-u/?}
  \end{itemize}
\end{encadre}

\begin{encadre}{Bande dessinés}
  \begin{itemize}
    \item \textit{Le monde sans fin}, Jancovici et Blain (2021)
    \item \textit{Ressources : un défi pour l'humanité}, Bihouix et Perriot (2024)
  \end{itemize}
\end{encadre}

\begin{encadre}{Vidéos}
  \begin{itemize}
    \item \textit{Climatisation : un désastre écologique ?}, \textit{Le Reveilleur}, R.Meyer (2026), \url{https://www.youtube.com/watch?v=nERo6YzNoSM}
  \end{itemize}
\end{encadre}